% Complete independent proof for the exact O/empty-base hierarchy.
\begin{annexitem}{annex:p04.doh.arithmetic-graphs}{thm:p04-doh-arithmetic-graphs}{Arithmetic formulas for computable numerical graphs}
\begin{proof}
We give an explicit finite-sequence resource first. Define
\(\beta(c,d,i)=r\) by the arithmetical remainder condition
\[
 r<1+(i+1)d\quad\text{and}\quad
 \exists q\,[c=q(1+(i+1)d)+r].
\]
Every finite sequence \(a_0,\ldots,a_{l-1}\) has such a code.
Choose positive \(d\) larger than all its entries and divisible by \(l!\).
The empty sequence imposes no remainder requirements. The moduli
\(m_i=1+(i+1)d\) are pairwise coprime. Indeed, for \(i<j<l\), a prime
dividing \(m_i,m_j\) cannot divide \(d\), but divides
\((j-i)d\), hence \(j-i\). It then divides \(l!\) and \(d\),
a contradiction. For coprime \(u,v\), the Euclidean algorithm gives
integers \(r,s\) with \(ru+sv=1\): successively divide with smaller
remainders until the last nonzero remainder and substitute backwards.
Consequently the two congruences modulo \(u,v\) can be solved by
\(b\,ru+a\,sv\), reduced modulo \(uv\), for prescribed residues
\(a,b\). Induction combines the finitely many pairwise coprime moduli.
The resulting nonnegative \(c\) satisfies \(\beta(c,d,i)=a_i\)
for every \(i<l\). All conditions on a proposed code use only
number quantifiers, addition, multiplication and order; order itself is
\(x<y\) iff \(\exists z\,y=x+(z+1)\).

Code tuples by the numerical pairing
\(2\langle x,y\rangle=(x+y)(x+y+1)+2y\), iterated for longer tuples.
It has an arithmetical graph and unique components: if
\(w=x+y\), its codes range from \(w(w+1)/2\) to
\((w+1)(w+2)/2-1\); these disjoint consecutive intervals cover
\(\N\). Within one interval \(y\), and then \(x\),
are uniquely recovered. Existence follows by finding the greatest
\(w\leq z\) whose triangular number is at most the proposed code \(z\).

For a fixed Turing machine, store a configuration as its state, scanned
symbol, and the finite tape words to the left and right of the head.
Encode a word as base-\(b\) digits, for a fixed alphabet size \(b\),
with omitted leading blanks understood as blanks. Reading or removing
a digit is a remainder/division condition, and adjoining a digit is
multiplication by \(b\) followed by addition. Each of the finitely
many machine instructions therefore gives an arithmetic relation
\(\operatorname{Step}(u,v)\) between coded configurations; their
finite disjunction is the transition relation. Initial input and output
conditions are arithmetic as well. If an input convention needs a word
of length \(n\), its digits can be specified by the sequence coding
just proved, or by a sequence satisfying the finite recurrence for the
required powers of \(b\).

There is now an arithmetic formula saying that a \(\beta\)-coded
sequence begins with the specified input configuration, each adjacent
pair satisfies \(\operatorname{Step}\), and its last configuration
halts with output \(y\). Quantify the sequence code and its length
existentially and check its finitely many transitions by a bounded
universal number quantifier. A genuine halting computation supplies
this finite code. Conversely any code satisfying these clauses describes
that computation, since its initial condition and every transition are
the machine's actual ones. This proves the graph equality, including
undefined inputs, which have no successful history. The construction
is effective from the finite instruction table. Predicates follow by
machines outputting their characteristic bits, and total functions have
their unique actual output in the standard model.
\end{proof}
\end{annexitem}

\begin{annexitem}{annex:p04.doh.program-tools}{thm:p04-doh-program-tools}{Program self-reference and uniform jump simulation}
\begin{proof}
Let \(s(p,t)\) insert the first numerical argument \(t\) into a program
with index \(p\). Thus
\(\varphi_{s(p,t)}^B(x)=\varphi_p^B(t,x)\), and \(s\) is total
computable independently of \(B\). For a given total computable \(f\),
let \(p\) index the program
\[
 (t,x)\longmapsto\varphi_{f(s(t,t))}^B(x).
\]
Put \(e=s(p,p)\). For every oracle \(B\),
\(\varphi_e^B(x)=\varphi_p^B(p,x)=\varphi_{f(e)}^B(x)\).
Tupling additional arguments proves the parameter version. The proof only
computes a finite program index; it makes no assertion that that program
terminates.

To read \(B(x)\) from \(B'\), effectively construct a program which
ignores its input, queries \(B(x)\), and halts exactly when that answer
is \(1\). Its index belongs to \(B'\) exactly when \(x\in B\).
This yields a fixed uniform reduction, and repeating it reads \(B\)
or \(B'\) from \(B''\).

If \(\varphi_d^B=\chi_A\), construct, from \(d,x\), a program
\(q_{d,x}\) with oracle \(B\) which ignores its own input and
simulates \(\varphi_x^A(x)\), answering each query \(A(u)\) by
running \(\varphi_d^B(u)\). All these subsidiary computations halt
and answer correctly. Hence
\[
 x\in A'\quad\Longleftrightarrow\quad
 \operatorname{index}(q_{d,x})\in B'.
\]
The index of this simulator is computable from \(d,x\), without knowing
\(A\) or \(B\). Oracle substitution and numerical specialization
are finite syntactic program operations; these are the required uniformities.
\end{proof}
\end{annexitem}

\begin{annexitem}{annex:p04.doh.notation-existence}{thm:p04-doh-notation-existence}{Denotations and the jump hierarchy}
\begin{proof}
The positive simultaneous rules act on a subset of the set consisting
of tagged natural numbers and tagged numerical pairs. Starting empty,
retain the old records and apply the rules at successors; take unions
at limits. This monotone process stabilizes: otherwise the distinct
stage sets indexed below the Hartogs ordinal of its powerset would
give an injection of that ordinal into the powerset. The stabilized
set is closed, and induction puts every stage below every closed set,
so it is the least simultaneous closure.

At the first admission of a notation, assign its denotation. For \(1\)
assign \(0\); for \(\mathsf s(a)\) assign the already assigned
denotation of \(a\) plus one. At the admission of \(\mathsf l(e)\),
all sequence terms have already been admitted and their successive
predecessor relations have been established. Assign the supremum of
their denotations. Simultaneously induct on record-admission stages
that every predecessor pair has strictly increasing denotations.
A successor pair follows from the preceding inequality and ordinal
successor. A limit predecessor lies at or below some sequence term;
that term is strictly below the next term, so its denotation is
strictly below the supremum. Previously admitted pairs retain their
inequalities. The same reasoning applies to further predecessor
pairs admitted after their upper notation first appeared.

The three numerical constructor forms are disjoint. A successor code
uniquely specifies its predecessor, and a limit code uniquely specifies
its program and therefore its total sequence. Consequently the
denotation never depends on the stage at which the same constructor
is reconsidered. Induction on its children proves uniqueness.
Every child of a valid notation has strictly smaller denotation.
No infinite child path is possible, since ordinals have no infinite
strict descending sequence.

For completeness the absence of a descending sequence gives ordinary
well-foundedness on this numerical domain without a choice axiom:
if a nonempty subset had no minimal element, start at its least
natural-number member and repeatedly take the least smaller-related
member in that subset. This gives the prohibited sequence.
Well-founded set recursion therefore assigns the sets \(H_a\) using
the three displayed clauses. Successor jump and numerical join are
set operations; Separation and Replacement collect their values.
Their uniqueness follows by the same well-founded induction.
\end{proof}
\end{annexitem}

\begin{annexitem}{annex:p04.doh.addition}{thm:p04-doh-addition}{A notation compiler and its uniform oracle reduction}
\begin{proof}
First construct the compiler. Given a proposed program index \(p\),
write a program taking \(a,b\) which returns \(a\) at \(b=1\),
recurses on the smaller code \(c\) at \(b=2^c>1\), and at
\(b=3\cdot5^e\) returns the limit code of the specialized program
\(i\mapsto\varphi_p(a,\varphi_e(i))\). On other codes it returns
\(1\). The self-reference lemma supplies an index whose behavior
obeys these clauses. The limit branch only constructs an index. A
top-level computation can recurse only through finitely many strictly
smaller natural numbers in the successor branch; its base, limit and
invalid branches terminate. Thus the compiler is total on all inputs.

Fix valid \(a\). Induct on the denotation of the second argument.
The base output is valid. At a successor the induction hypothesis
gives a valid predecessor and the denotation adds one. At a limit
\(b=\mathsf l(e)\), the sequence
\(\operatorname{add}_{\mathcal O}(a,\varphi_e(i))\) is valid and
strictly increasing in notation order, as established below for its
earlier terms. Its limit code is therefore valid. Continuity of
ordinal addition in its second argument gives the asserted denotation.

Include predecessor preservation in this induction. If
\(x<_{\mathcal O}\mathsf s(c)\), then \(x=c\) or
\(x<_{\mathcal O}c\). The compiled successor puts the compiled
\(c\) immediately below it; induction and transitivity handle the
second case. If \(x<_{\mathcal O}\mathsf l(e)\), then
\(x\leq_{\mathcal O}\varphi_e(i)\) for some \(i\).
Its compiled value is at or below that compiled sequence term and
therefore below the compiled limit. This proves preservation,
including the successive comparisons needed in the preceding limit
case. Also include \(a\leq_{\mathcal O}\operatorname{add}_{\mathcal O}(a,b)\)
in the induction. It is equality at \(b=1\). At a successor the
outer successor makes it strict. At a limit a strictly increasing
sequence has a nonzero term; the already proved strict inequality
for that term followed by the limit inequality makes it strict.
Transitivity here is part of the predecessor clauses: each upper
notation contains all predecessors of its selected earlier terms.

Now construct the right-summand reduction, rather than invoking an
ordinal comparison theorem. Fix by self-reference a universal oracle
program \(M(a,b,x)\) with the following behavior; specialization
gives the total index function \(r(a,b)\). Its intended oracle is
\(Y=H_{\operatorname{add}_{\mathcal O}(a,b)}\).
At \(b=1\) it returns \(0\), deciding \(H_1=\varnothing\).
At \(b=\mathsf s(c)\), specialize the same program to \(a,c\),
and apply the jump simulator of the first lemma to that reduction.
The resulting computable index \(q_x\) is a query to \(Y\), because
\[
 Y=H_{\operatorname{add}_{\mathcal O}(a,c)}'.
\]
Return the oracle bit \(Y(q_x)\). The specialized child index is
computed syntactically, without running the reduction on all inputs.

At \(b=\mathsf l(e)\), decode \(x=\langle i,u\rangle\), compute
\(c=\varphi_e(i)\), and run \(M(a,c,u)\) with each oracle query
\(v\) answered by \(Y(\langle i,v\rangle)\). This is exactly its
required oracle, since the compiled limit sequence is
\(i\mapsto\operatorname{add}_{\mathcal O}(a,\varphi_e(i))\).
Invalid constructor cases can return \(0\); no correctness claim
is made for them. These operations specify an effective program
transformation, so the self-reference lemma supplies the one index.

Finally induct on \(\nu_{\mathcal O}(b)\). The base case is correct
and total. The successor jump simulator is correct by the already
total child reduction. At a valid limit \(\varphi_e(i)\) halts,
and the chosen child has smaller denotation, so its reduction halts
and gives the correct join membership. This proves both correctness
and totality on every valid input and completes the required uniform
reduction. It does not decide membership in \(\mathcal O\).
\end{proof}
\end{annexitem}

\begin{annexitem}{annex:p04.doh.node-bounds}{thm:p04-doh-node-bounds}{Effective valid bounds and explicit child columns}
\begin{proof}
Given a prospective generator index \(p\), compile a sequence program
which, with \(z\) inserted as a parameter, calculates \(a_{z,i}\)
by its finite loop. In a present-child case that loop calls
\(\varphi_p(i)\); in an absent-child case it uses \(1\). Define
the transformed generator's output at \(z\) to be the two successive
successor wrappers around the limit code of this specialized sequence
program. This output computation does not call \(\varphi_p\).
Self-reference fixes \(p\), and its top-level behavior is therefore
a total computable function \(h(z)\). Since notation addition is
also total, every individual sequence value is total computable.

Use well-founded induction on the child relation. Suppose every
child's \(h(i)\) is valid. In the absent-child case the dummy \(1\)
is valid. Starting at \(1\), finite induction and notation addition
prove that every \(a_{z,i}\) is valid and
\[
 a_{z,i}<_{\mathcal O}\mathsf s(a_{z,i})
       \leq_{\mathcal O}a_{z,i+1}.
\]
The sequence thus satisfies the actual \(\mathcal O\) limit
condition. Its limit and its two successors are valid, proving
\(h(z)\in\mathcal O\). Well-founded induction applies on the
natural-number domain by the least-child argument in the preceding
notation proof; no selection of arbitrary branches is involved.

For a present child \(i\), column \(i+1\) of \(J_z\) is
\[
 H_{a_{z,i+1}}=
 H_{\operatorname{add}_{\mathcal O}(\mathsf s(a_{z,i}),h(i))}.
\]
The previous lemma's specialized reduction reads \(H_{h(i)}\)
from exactly this column. Both \(a_{z,i}\) and the reduction index
are computable from \(z,i\); oracle-column substitution is a finite
program operation. This proves uniform child access from \(J_z\).
The two jump-downgrade simulators read \(J_z\) from \(J_z''\),
proving the final assertion without an assumed general effective
transfinite-recursion closure theorem.
\end{proof}
\end{annexitem}

\begin{annexitem}{annex:p04.doh.alternation}{thm:p04-doh-alternation}{The two actual jump queries and the separating bit}
\begin{proof}
Well-founded recursion on the synchronized product defines a unique
bit at each pair; its clause only uses immediate child pairs.
Here is its explicit oracle program. Fix a prospective evaluation
index \(d\). At a pair node \(z=(n,s,t)\), use oracle \(J_z\)
to compute the following predicate in \(m,k\):
\[
 P_z(m,k)=
 \bigl[s^\frown m\notin T_n\bigr]\ \lor\
 \bigl[t^\frown k\in U_n\ \land\
     \varphi_d^{H_{h(c)}}(c)=1\bigr],
 \qquad c=(n,s^\frown m,t^\frown k).
\]
The second disjunct is only run if the left child is present and
the right child is present. Thus \(c\) is then an actual product
child. Its oracle \(H_{h(c)}\) is simulated from \(J_z\) by
the exact column reduction in the preceding lemma. Both tree tests
are computable. Under the inductive promise that child evaluations
are total, this predicate is total relative to \(J_z\).

For each \(m\), compile a program \(p_{z,m}\) with oracle \(J_z\)
which ignores its own input, searches \(k=0,1,\ldots\), and halts
as soon as \(P_z(m,k)=1\). Consequently
\[
 \operatorname{index}(p_{z,m})\in J_z'
       \quad\Longleftrightarrow\quad\exists k\,P_z(m,k).
\]
Now compile a program \(q_z\) with oracle \(J_z'\) which ignores
its own input, searches \(m=0,1,\ldots\), and halts when its oracle
answers \(0\) at \(\operatorname{index}(p_{z,m})\). Then
\[
 \operatorname{index}(q_z)\in J_z''
       \quad\Longleftrightarrow\quad
       \exists m\,\neg\exists k\,P_z(m,k).
\]
The transformed evaluation program, whose oracle is
\(Y=H_{h(z)}=J_z''\), returns
\(1-Y(\operatorname{index}(q_z))\). Constructing \(q_z\)
only constructs finite indices containing \(d,z\); it does not
run all child evaluations. On numbers not coding pair nodes let
the evaluator return \(0\). Program self-reference supplies the
one index \(d\) whose behavior is this transformation.

Induct on the well-founded product-child relation. All child evaluations
are correct by induction, so \(P_z\) is total and has exactly the
displayed predicate values. The two halting equivalences are therefore
valid, and the returned bit is exactly \(\forall m\exists k\,P_z(m,k)\),
which is the defining equation. The evaluator itself terminates after
constructing indices and reading one bit of its stipulated oracle.
This proves correctness and totality at every node. In particular
leaf, empty-existential and empty-universal cases are included, not
left to an unspecified stopping rule.

Let \(i_n\) be the numerical code of the root pair of the \(n\)-th
product tree. It is a child of \(*\). Using column \(i_n+1\) of
\(J_*\) and its specialized notation-addition reduction, simulate
\(H_{h(i_n)}\), then run the just constructed evaluator at \(i_n\).
Two jump downgrades supply \(J_*\) from \(H_{h(*)}=J_*''\).
All indices and \(i_n\) are computable from \(n\). This is one
total decision program with one fixed valid oracle, not a choice
of a different unspecified ordinal bound for each input.

For the separation property, suppose \(U\) has a branch. Induct
on a well-founded subtree of \(T\) that a pair with a node on
this branch has bit \(1\). For each \(T\)-child, use the next
child on the branch of \(U\); induction gives its bit \(1\).
The universal clause holds, including a leaf of \(T\).
Conversely suppose \(T\) has a branch and \(U\) is well founded.
Induct on \(U\) that a pair whose first node lies on the branch
has bit \(0\). Select the next \(T\)-child on that branch.
Every available \(U\)-child pair has bit \(0\) by induction,
so the existential clause fails for this selected first child.
At a leaf of \(U\) it fails because there is no second child.
This proves the reversed case. Each argument uses one given branch,
not an effective procedure for selecting branches.
\end{proof}
\end{annexitem}

\begin{annexitem}{annex:p04.doh.tree-normal-form}{thm:p04-doh-tree-normal-form}{Recursive witness trees}
\begin{proof}
First match the actual function-quantifier convention to set quantification.
For a set \(X\subseteq\N\), the arithmetic guard
\[
 \forall i\,\exists j\,
 \bigl[\langle i,j\rangle\in X\ \land\
       \forall k\,(\langle i,k\rangle\in X\longrightarrow k=j)\bigr]
\]
says exactly that \(X\) is the graph of a total unary numerical function.
The pairing is the arithmetic pairing already constructed. Every
\(\alpha\in\N^{\N}\) has such a graph by Separation, and a guarded
graph determines exactly one \(\alpha\) by its unique value at every
input. Replacement collects these values without a choice principle.

Translate a formula with \(\alpha\) by introducing a fresh numerical
variable for the value of each function-application subterm. At an
application to an input with value \(i\), require the output value \(j\)
to satisfy \(\langle i,j\rangle\in X\). Arithmetic constructors use
their ordinary arithmetic graphs. At an atomic formula existentially
quantify its finitely many intermediate numerical values and impose the
original relation on its resulting term values. Then preserve Boolean
connectives and number quantifiers, renaming the fresh variables to avoid
capture. Under the guard, term induction gives exactly one value for
every subterm, and formula induction proves agreement with the original
formula, including negated atoms and nested function applications.
The resulting matrix is arithmetical in \(X\): the extra quantifiers
are all number quantifiers.

Thus replace an existential function quantifier by an existential set
quantifier with this guard conjoined to the translated matrix. Existence
of an actual function supplies its guarded graph in one direction; a
guarded graph supplies the unique actual function in the other. Replace
a universal function quantifier by a universal set quantifier with the
guard implying the translated matrix. Invalid graph sets then impose no
condition, while every actual function is represented and every valid
graph represents an actual function. This proves both directions of
the universal equivalence, not just preservation of existential witnesses.

Conversely, for a set variable \(X\) use a function \(\alpha\) subject
to the arithmetic binary-value guard
\[
 \forall i\,[\alpha(i)=0\ \lor\ \alpha(i)=1],
\]
and replace each atom \(t\in X\) by \(\alpha(t)=1\). Every set
has its unique binary characteristic function, and every guarded function
represents the set of inputs on which it is \(1\). Formula induction
again proves agreement, including negations and number quantifiers.
Conjoin the binary guard for an existential function quantifier and use
it as an implication antecedent for a universal function quantifier.
This proves both existential and universal set-to-function equivalences.

All these transformations are effective operations on a finite formula
and add only arithmetic number quantifiers to its matrix. For finitely
many variables use one guard per variable and pack their graphs or
functions through the fixed pairing. No extra second-order quantifier
alternation is introduced. In particular the lightface \(\Sigma^1_1\)
and \(\Pi^1_1\) classes in the original function convention are
exactly the classes used in the set-based proof, with no set parameter
introduced. Applying this translation, continue with a universal-set
definition for the tree construction.

Negate the given universal-set definition. Consider its form
\(\exists X\,\eta(n,X)\), with an arithmetical matrix. Represent
\(X\) by its binary characteristic function. Put the finitely
many number quantifiers in prenex form, after capture-free renaming.
For each existential number quantifier use a witness function of its
preceding universally quantified number variables. If the prefix is
true for a given \(X\), these functions exist: at each tuple of
preceding universal numbers choose the least natural witness leaving
the remaining suffix true. Induction on the finite quantifier prefix
verifies this selection and the converse implication. It is numerical
least-witness selection, not countable choice for arbitrary families.

Pack the binary characteristic function and these finitely many witness
functions into one numerical function \(g\), using pairing to specify
the component and its argument tuple. It must satisfy a recursive
list of tests: the characteristic component is binary, and the
quantifier-free matrix holds at every tuple of universal arguments.
Every test is a halting oracle computation using finitely many entries
of \(g\). Form a recursive tree by rejecting a finite sequence
\(s\) when one of the first \(\operatorname{len}(s)\) tests,
run for \(\operatorname{len}(s)\) steps using its entries, finishes
and reports failure. A missing entry leaves a test unfinished.
Include the empty sequence. A rejected prefix remains rejected on
an extension: the same failing test is included, gets at least as
much time, and receives the same earlier entries. Thus the accepted
sequences are prefix closed.

A successful \(g\) has all prefixes in this tree. Conversely a
branch cannot fail a test: a failure would use finitely many entries
and steps, and a sufficiently long prefix would be rejected. Hence
the tree has a branch exactly when the existential-set assertion
holds. The entire program is compiled from the finite formula and
the numerical parameter \(n\). Negating its branch condition
gives the stated well-founded-tree normal form for the original
\(\Pi^1_1\) definition. Computable transition and coding operations
have arithmetic graphs via their finite numerical computation histories;
no set parameter or hierarchy theorem enters the construction.
\end{proof}
\end{annexitem}

\begin{annexitem}{annex:p04.doh.delta-to-hyp}{thm:p04-doh-delta-to-hyp}{One oracle for the complementary pair}
\begin{proof}
Apply the tree-normal-form compiler to the two given definitions,
getting recursive trees \(T_n,U_n\) satisfying
\[
 n\in A\ \Longleftrightarrow\ T_n\text{ is well founded},\qquad
 n\notin A\ \Longleftrightarrow\ U_n\text{ is well founded}.
\]
For each \(n\), exactly one tree is well founded and the other
has a branch. The synchronized product has a branch only if both
trees have branches: its two projections would be such branches.
It is therefore well founded. Conversely, on numerical tree domains
absence of a branch is ordinary well-foundedness, by the least-child
argument already given. Add a single common node with all these
root pairs as children. An infinite child sequence starting at this
node would choose one index \(n\) and thereafter remain in that
one well-founded product, which is impossible. Nodes with invalid
codes have no children. The whole numerical child relation is
recursive and well founded.

Compile its node generator \(h\). The preceding lemmas give the
valid code \(a=h(*)\) and the total program for all product-root
bits using oracle \(H_a\). If \(n\in A\), the first tree is
well founded and the second has a branch, so its bit is \(1\).
If \(n\notin A\), the reversed case gives bit \(0\). Thus
this one program decides \(A\) from \(H_a\).

Every program, node code and notation code used here is compiled
from the two original finite definitions. The mathematical proof
uses their complementary promise to prove validity and correctness;
the compiler does not test it. Since a \(\Delta^1_1\) set and its
complement both have \(\Pi^1_1\) definitions, this proves the
claimed inclusion in the already defined iterated-jump class.
\end{proof}
\end{annexitem}

\begin{annexitem}{annex:p04.doh.hyp-to-delta}{thm:p04-doh-hyp-to-delta}{The unique arithmetically described hierarchy}
\begin{proof}
Fix a valid notation \(a\). Unfold its constructors into a tree
\(V_a\) of finite choice sequences, with the empty node labelled
by \(a\). A node labelled \(1\) has no children; a node labelled
\(\mathsf s(b)\) has one child, numbered \(0\), labelled \(b\);
a node labelled \(\mathsf l(e)\) has child \(i\) labelled
\(\varphi_e(i)\) for each natural \(i\). Membership and labels
are computed by following the finite choices. This halts for every
finite sequence when \(a\) is valid: every reached limit sequence
is total, and each reached notation is valid. At the first invalid
choice reject the sequence. The tree is therefore recursive for
this fixed valid \(a\). Uniformly for arbitrary numerical \(a\),
its membership and label computations have arithmetical graph
definitions; no uniform totality on invalid codes is required.
All child labels have strictly smaller ordinal denotation, so the
tree is well founded.

Let \(Y\subseteq\N\) code columns \(Y_u\subseteq\N\), one for
each finite-sequence node \(u\), by numerical pairing. Write
\(\operatorname{Hier}(a,Y)\) for the following requirements:
columns outside \(V_a\) are empty; a base-labelled column is empty;
at a successor-labelled node \(u\),
\[
 x\in Y_u\quad\Longleftrightarrow\quad
 \varphi_x^{Y_{u^\frown0}}(x)\text{ halts};
\]
and at a limit-labelled node \(u\),
\[
 \langle i,x\rangle\in Y_u\quad\Longleftrightarrow\quad
 x\in Y_{u^\frown i}.
\]
This is one arithmetical relation in the set predicate \(Y\) and
the numerical parameter \(a\). Universal quantification over node
codes and column entries uses only number quantifiers. Halting with
a column oracle is stated by existence of a finite computation
history whose every oracle-answer bit agrees with that column of
\(Y\). The history, indices, lengths and transitions are numerical;
their checks are computable, hence arithmetically expressible. There
is no hidden quantifier over sets of histories or over truth values.

Well-founded set recursion on \(V_a\) constructs a solution \(Y\),
using its empty, jump and join clauses; Replacement collects the
countably many columns and Separation forms their paired union.
The same induction proves uniqueness. A second induction, comparing
these clauses with the definition of \(H\), proves that each
column equals the set for its node's notation. In particular
\(Y_{\varnothing}=H_a\). Repeated occurrences of the same notation
give the same set by uniqueness of the notation hierarchy; they do
not supply competing choices of column values.

Consequently
\[
\begin{aligned}
 x\in H_a&\quad\Longleftrightarrow\quad
   \exists Y\,[\operatorname{Hier}(a,Y)\land x\in Y_{\varnothing}],\\
 x\notin H_a&\quad\Longleftrightarrow\quad
   \exists Y\,[\operatorname{Hier}(a,Y)\land x\notin Y_{\varnothing}].
\end{aligned}
\]
Existence ensures that both formulas give an answer; uniqueness
ensures that they give complementary answers. Both are lightface
\(\Sigma^1_1\) with the fixed natural \(a\) inserted as a numeral.

If \(\varphi_d^{H_a}=\chi_A\), replace their last matrix condition
by existence of a finite computation of \(\varphi_d^{Y_{\varnothing}}(x)\)
halting with output \(1\), respectively output \(0\). The total
reduction and the unique existing hierarchy again make these two
existential-set arithmetic definitions complementary. Thus every
set computable in \(H_a\) is \(\Delta^1_1\). The same fixed
finite formulas accept \(a,d\) as number parameters, proving the
stated uniformity on valid, total inputs.
\end{proof}
\end{annexitem}

\begin{annexitem}{annex:p04.doh.characterization}{thm:p04-doh-characterization}{Equality and the parameter boundary}
\begin{proof}
The two independently proved inclusions give the equality for
subsets of \(\N\). Numerical pairing transports it to finite
arity relations. Inserting a fixed finite tuple of numbers into a
formula or program is effective, so numerical parameters do not
change the pointclass or the computation class.

The equality over arithmetic uses no arbitrary set parameter. In
particular, all limit sequence programs in the proved statement are
computable without an oracle. Changing those programs to computations
relative to a set is a change of the notation system, not insertion
of an additional natural numeral into the preceding construction.
All well-founded recursions and numerical least selections used in
the proof are available in classical ZF; arbitrary countable choice
is not a dependency.
\end{proof}
\end{annexitem}

\begin{annexitem}{annex:p04.doh.overview}{thm:p04-hyp-delta-one-one}{The existing hyperarithmetical characterization}
\begin{proof}
Definition~\ref{def:p04-doh-notations} supplies exactly the simultaneous
\(\mathcal O\) convention used by
Definition~\ref{def:p04-kleene-ordinal-notations}, with empty base jump set,
successor jump and effective limit join. The construction and existence
proof retain the notation-code order rather than replacing it by the order
of denoted ordinals. The effective function/set-quantifier bridge in
Lemma~\ref{thm:p04-doh-tree-normal-form} identifies the same old lightface
pointclasses. The two proved inclusions in
Theorem~\ref{thm:p04-doh-delta-to-hyp} and
Theorem~\ref{thm:p04-doh-hyp-to-delta} therefore establish the stated
equality, as summarized by Corollary~\ref{thm:p04-doh-characterization}.
\end{proof}
\end{annexitem}
